\clearpage
\thispagestyle{empty}
\pagestyle{empty}
\null
\begin{tikzpicture}[remember picture,overlay,x=1in,y=1in,shift={(current page.south west)}]
  \MarkCoverGround
  \node[text=MarkCoverMuted,font=\fontsize{11}{14}\selectfont] at (3.5,9.25)
    {Υἱὸς Θεοῦ \quad 神的兒子};
  \node[text=MarkCoverGold,font=\fontsize{22}{28}\selectfont\bfseries] at (3.5,8.82)
    {先做兒子，才能作僕人};
  \node[text=MarkCoverMuted,font=\fontsize{13}{17}\selectfont\itshape] at (3.5,8.46)
    {Son First, Then Servant};
  \draw[MarkCoverGold,line width=0.5pt] (2.80,8.05) -- (4.20,8.05);
  \node[anchor=north,text=MarkCoverCream,align=left,text width=5.15in,
        font=\fontsize{12}{19}\selectfont] at (3.5,7.66)
    {從加利利的服事，到耶路撒冷的十字架，馬可福音讓我們看見：神的兒子，如何走完僕人之路。\\[0.17in]
     本書逐章研讀馬可福音十六章，以 1:1 的身分宣告為起點、10:45 的捨命贖價為心臟，整合老弟兄查經法、John MacArthur 逐節解經與 G. Campbell Morgan 的屬靈組織分析，並貫通四福音、以賽亞書與希伯來書。};
  \draw[MarkCoverGold,opacity=0.6,line width=0.4pt] (0.93,5.58) -- (6.07,5.58);
  \foreach \x/\n/\cn/\en in {1.75/I/經文與原文/Scripture,3.5/II/結構與注疏/Exposition,5.25/III/領受與操練/Formation}{%
    \node[text=MarkCoverGold,font=\fontsize{11}{14}\selectfont] at (\x,5.25) {\n};
    \node[text=MarkCoverCream,font=\fontsize{11}{15}\selectfont] at (\x,4.94) {\cn};
    \node[text=MarkCoverMuted,font=\fontsize{10}{13}\selectfont\itshape] at (\x,4.67) {\en};
  }
  \draw[MarkCoverGold,opacity=0.6,line width=0.4pt] (0.93,4.36) -- (6.07,4.36);
  \node[anchor=north,text=MarkCoverMuted,align=left,text width=5.10in,
        font=\fontsize{11}{16}\selectfont] at (3.5,4.02)
    {A bilingual study of all sixteen chapters of Mark, tracing the Son of God from His ministry in Galilee to His sacrifice and resurrection. Scripture, word studies, commentary and reflection invite the reader to move from understanding the text to a life of faithful service.};
  \node[text=MarkCoverCream,align=center,text width=5.2in,
        font=\fontsize{12}{19}\selectfont] at (3.5,2.43)
    {「因為人子來，並不是要受人的服事，乃是要服事人，\\並且要捨命作多人的贖價。」\\[0.11in]
     {\fontsize{9}{12}\selectfont\color{MarkCoverGold}馬可福音 10:45 \textperiodcentered\ Mark 10:45}};
  \draw[MarkCoverGold,line width=0.5pt] (2.95,1.65) -- (4.05,1.65);
  \node[text=MarkCoverCream,align=center,font=\fontsize{10}{14}\selectfont] at (3.5,1.20)
    {馬可福音研讀 \textperiodcentered\ 2026 整編版\\[0.08in]
     {\color{MarkCoverMuted}PubHub 三書精讀系統 \textperiodcentered\ 三書精讀出版系統}\\[0.06in]
     {\color{MarkCoverGold}Soli Deo Gloria}};
\end{tikzpicture}
